\documentclass[11pt,a4paper]{article}
\usepackage[flat]{pinns-course}

\coursemeta{Session 1, problem class --- Universal approximation}{Session 1, problem class}{universal approximation, Cybenko, Leshno-Lin-Pinkus-Schocken, neural networks, M2 MACIA, AGM, CY Cergy Paris University}

% ================================================================
\begin{document}

\coursetitle{Session 1 \textbullet\ Problem class}{Universal approximation}{Rebuilding the theorems of Cybenko and of\\[2pt] Leshno--Lin--Pinkus--Schocken from scratch}

\vspace{6pt}

\begin{plainbox}
\textbf{How this sheet works.} It replaces the lecture on universal approximation. Everything that
would have been proved at the board is proved here by you. Each part opens with a short box giving
the objects and the external results you may use \emph{without proof}; everything else is to be
established.

\smallskip
Exercises marked \textsc{core} carry the main line of argument and should be done in order: together
they reconstruct the two central theorems. Those marked \textsc{if time permits} probe the
hypotheses and can be taken home.

\smallskip
Hints are collected in §\,8, deliberately away from the statements. Full solutions are in §\,9.

\smallskip
\textbf{Prerequisites.} Session 1, §\,1 (the class $\Sigma_n(\sigma)$ and its structure).
Functional analysis: Hahn--Banach, dual of $C(K)$, Riesz--Markov. Measure theory: dominated
convergence, Fourier transform of a finite measure. Stone--Weierstrass.
\end{plainbox}

\newpage
\tableofcontents
\newpage

% ================================================================
\section{Notation, and the question}

\begin{setting}
Throughout, $d \ge 1$, $K \subset \R^d$ is compact, and $\sigma : \R \to \R$ is an activation.
Recall from §\,1 of the lecture notes:
\[
\Sigma_n(\sigma) := \Big\{ x \mapsto \sum_{i=1}^n c_i\, \sigma(w_i \cdot x + b_i)
\ :\ c_i \in \R,\ w_i \in \R^d,\ b_i \in \R \Big\},
\qquad
\Sigma(\sigma) := \bigcup_{n \ge 1} \Sigma_n(\sigma),
\]
so that $\Sigma(\sigma)$ is the linear span of the \kw{ridge functions}
$x \mapsto \sigma(w\cdot x + b)$, and is a one-hidden-layer object. $C(K)$ carries the uniform norm
$\norm{f}_\infty = \sup_K\abs f$.
\end{setting}

\noindent
\textbf{The question of this sheet.} For which $\sigma$ is $\Sigma(\sigma)$ dense in $C(K)$?

\smallskip
\noindent
Two answers will be reconstructed. The first (Part D--E) is Cybenko's 1989 theorem: density holds
for every continuous bounded \emph{sigmoidal} $\sigma$. The second (Part F) is the complete
characterisation of Leshno, Lin, Pinkus and Schocken: for continuous $\sigma$, density holds
\emph{if and only if $\sigma$ is not a polynomial}. Part B proves the ``only if'', which is easy and
instructive; Part G asks what each hypothesis is really doing.

\begin{warnbox}
\textbf{A word on what we are \emph{not} doing.} Nothing on this sheet produces a network, an
algorithm, or a bound on the number of neurons needed to reach accuracy $\varepsilon$. Keep that in
view: Exercise \ref{ex:reading} is where we make the observation explicit, and it is the reason
§\,3 of the lecture notes exists at all.
\end{warnbox}

% ================================================================
\section{Part A --- Getting the object in hand}

\begin{setting}
A function $\sigma$ is \kw{sigmoidal} if it is bounded and measurable with
\[
\lim_{t\to+\infty}\sigma(t) = 1, \qquad \lim_{t\to-\infty}\sigma(t) = 0 .
\]
The logistic function $\mathrm{sig}(t) = (1+\ee^{-t})^{-1}$ is sigmoidal and continuous; so is
$\tfrac12(1+\tanh(t/2))$. No result is needed for this part beyond the definitions.
\end{setting}

\begin{exercise}[Warm-up: what does $\Sigma(\sigma)$ contain? \tagcore]\label{ex:warmup}
\leavevmode
\begin{enumerate}[label=\textbf{\arabic*.}]
\item Take $d = 1$ and $\sigma = \tanh$. Describe in words what each of $c$, $w$, $b$ does to the
graph of $x \mapsto c\,\sigma(wx + b)$. Which of the three controls the \emph{steepness}?
\item Show that if $\sigma(t_0) \ne 0$ for some $t_0$, then every constant function belongs to
$\Sigma_1(\sigma)$.
\item Let $\sigma$ be sigmoidal and let $a < b$ in $\R$. Show that, for $d = 1$,
\[
\sigma\big(\lambda(x-a)\big) - \sigma\big(\lambda(x-b)\big) \xrightarrow[\lambda\to+\infty]{}
\ind_{(a,b)}(x) \qquad \text{for every } x \notin \{a,b\}.
\]
Which element of $\Sigma_2(\sigma)$ is on the left?
\item Deduce an informal argument that $\Sigma(\sigma)$ should be dense in $C([0,1])$. Which
classical approximation result are you implicitly invoking?
\item Now let $d \ge 2$. Explain in two or three sentences why the argument of item 4 does not
transpose, and what goes wrong. \emph{This is the reason Parts C--E are needed at all.}
\end{enumerate}
\end{exercise}

% ================================================================
\section{Part B --- The easy direction: polynomial activations}

\begin{setting}
$\Pp_m(\R^d)$ denotes the space of real polynomials in $d$ variables of total degree $\le m$; it is
finite-dimensional. You may use without proof that a finite-dimensional subspace of a normed vector
space is closed.
\end{setting}

\begin{exercise}[Polynomial activations are not universal \tagcore]\label{ex:poly}
\leavevmode
\begin{enumerate}[label=\textbf{\arabic*.}]
\item Let $\sigma$ be a polynomial of degree $m$. Show that for all $w \in \R^d$ and $b \in \R$ the
function $x \mapsto \sigma(w\cdot x + b)$ belongs to $\Pp_m(\R^d)$. Deduce
$\Sigma(\sigma) \subseteq \Pp_m(\R^d)$.
\item Deduce that $\overline{\Sigma(\sigma)} \subseteq \Pp_m(\R^d)$, and conclude that
$\Sigma(\sigma)$ is not dense in $C(K)$ as soon as $K$ has nonempty interior. Exhibit explicitly a
function of $C(K)$ that is not in the closure.
\item Where exactly did you use that $K$ has nonempty interior? (Exercise \ref{ex:hypotheses} comes
back to this.)
\item Take $\sigma(t) = t^2$ and $d = 1$. Determine $\Sigma(\sigma)$ \emph{exactly}. Is the
inclusion of item 1 an equality here?
\item Still with $\sigma(t) = t^2$: show that a network of \emph{fixed} depth $L$ and arbitrary
width realises only polynomials of degree $\le 2^{L-1}$, so that this class is not dense either.
\item Now let the depth vary. Using $xy = \tfrac14\big((x+y)^2 - (x-y)^2\big)$, show that
$\bigcup_{L,W}\Nn^{t^2}_{L,W}(d)$ contains \emph{every} polynomial, hence is dense in $C(K)$.
\item Reconcile items 5 and 6, and answer: over which class of functions does the
Leshno--Lin--Pinkus--Schocken theorem quantify? Why does the answer to that question change the
truth value of the statement?
\end{enumerate}
\end{exercise}

% ================================================================
\section{Part C --- The functional-analytic skeleton}

\begin{setting}
Two results from the functional analysis course, usable without proof.

\smallskip
\textbf{(HB)} \emph{Hahn--Banach, analytic form.} A continuous linear functional defined on a
subspace of a normed space $X$ extends to $X$ with the same norm.

\smallskip
\textbf{(RM)} \emph{Riesz--Markov.} For $K$ compact, every $L \in C(K)'$ is of the form
$L(f) = \int_K f\,\dd\mu$ for a unique finite signed regular Borel measure $\mu$ on $K$, and
$\norm{L}_{C(K)'} = \abs\mu(K)$. In particular $L = 0$ if and only if $\mu = 0$.
\end{setting}

\begin{definition}[Discriminatory]\label{def:discr}
Let $K \subset \R^d$ be compact. A bounded measurable $\sigma : \R\to\R$ is \kw{discriminatory on
$K$} if the only finite signed regular Borel measure $\mu$ on $K$ satisfying
\[
\int_K \sigma(w\cdot x + b)\,\dd\mu(x) = 0 \qquad \text{for all } w\in\R^d,\ b \in \R
\]
is $\mu = 0$.
\end{definition}

\begin{exercise}[From a closed proper subspace to a measure \tagcore]\label{ex:skeleton}
\leavevmode
\begin{enumerate}[label=\textbf{\arabic*.}]
\item Let $X$ be a normed space and $V \subsetneq X$ a \emph{closed} subspace. Prove: there exists
$L \in X'$ with $L \ne 0$ and $L|_V = 0$.
\emph{(Take $x_0 \notin V$, set $\delta := \dist(x_0,V) > 0$, define $L$ on $V \oplus \R x_0$ by
$L(v + t x_0) := t\,\delta$, check $\norm L \le 1$ and $L(x_0) = \delta$, then apply \textup{(HB)}.)}
\item Why is it essential that $V$ be closed? Give a one-line reason, and say where in item 1 it is
used.
\item Let now $\sigma \in C(\R)$ be bounded, and suppose $\sigma$ is discriminatory on $K$ in the
sense of Definition \ref{def:discr}. Prove that $\Sigma(\sigma)$ is dense in $C(K)$.
\emph{You are assembling items 1 and \textup{(RM)}; the whole proof is six lines.}
\item At which precise point of your argument did you use the continuity of $\sigma$?
\item Suppose someone gives you Definition \ref{def:discr} only for $K = [0,1]^d$. What extra step
would you need to obtain item 3 for a general compact $K$? (Do it.) Conclude that stating the
definition on an arbitrary compact is not a matter of taste.
\end{enumerate}
\end{exercise}

% ================================================================
\section{Part D --- The heart: sigmoidal implies discriminatory}

\begin{setting}
For $w \in \R^d$ and $b \in \R$ write
\[
H^+_{w,b} := \{x \in K : w\cdot x + b > 0\}, \qquad
\Pi_{w,b} := \{x \in K : w\cdot x + b = 0\} .
\]
You may use: the dominated convergence theorem for a finite signed measure; the fact that a
continuous function on a compact interval is a uniform limit of step functions; and the
\emph{uniqueness theorem for the Fourier--Stieltjes transform}: if $\mu$ is a finite signed Borel
measure on $\R^d$ with $\int \ee^{\ii w\cdot x}\dd\mu(x) = 0$ for every $w \in \R^d$, then $\mu = 0$.
\end{setting}

\begin{exercise}[Every bounded sigmoidal function is discriminatory \tagcore]\label{ex:discr}
Let $\sigma$ be bounded, measurable and sigmoidal, let $K \subset \R^d$ be compact, and let $\mu$ be
a finite signed regular Borel measure on $K$ with
\begin{equation}\label{eq:hyp}
\int_K \sigma(w\cdot x + b)\,\dd\mu(x) = 0 \qquad \text{for all } w, b .
\end{equation}
The goal is to prove $\mu = 0$.

\medskip
\noindent\emph{Step 1: collapsing the sigmoid to a step.}
\begin{enumerate}[label=\textbf{\arabic*.}]
\item Fix $w, b$ and $\varphi \in \R$, and set $\sigma_\lambda(x) := \sigma\big(\lambda(w\cdot x+b) + \varphi\big)$
for $\lambda > 0$. Explain why \eqref{eq:hyp} gives $\int_K \sigma_\lambda \dd\mu = 0$ for every
$\lambda > 0$. \emph{(Which $w$ and $b$ are you feeding to the hypothesis?)}
\item Compute $\lim_{\lambda\to+\infty}\sigma_\lambda(x)$, distinguishing three cases according to
the sign of $w\cdot x + b$. Note carefully what happens on $\Pi_{w,b}$.
\item Justify the exchange of limit and integral. Which domination do you use, and where does the
finiteness of $\mu$ enter?
\item Deduce that
\[
\mu\big(H^+_{w,b}\big) + \sigma(\varphi)\,\mu\big(\Pi_{w,b}\big) = 0 \qquad\text{for every }\varphi\in\R,
\]
and then, by letting $\varphi \to \pm\infty$, that $\mu(H^+_{w,b}) = \mu(\Pi_{w,b}) = 0$ for all
$w,b$.
\end{enumerate}

\medskip
\noindent\emph{Step 2: from half-spaces to the Fourier transform.}
\begin{enumerate}[label=\textbf{\arabic*.}, resume]
\item Fix $w \in \R^d$. Show that $\{w\cdot x : x \in K\}$ is \emph{contained in} a compact interval
$[\alpha,\beta]$, and give $\alpha$ and $\beta$. Why did I write ``contained in'' rather than
``equal to''? Give an example.
\item For $h : [\alpha,\beta]\to\R$ bounded measurable, set $F(h) := \int_K h(w\cdot x)\,\dd\mu(x)$.
Show that $F$ is linear and that $\abs{F(h)} \le \norm h_\infty \abs\mu(K)$.
\item Using Step 1, show that $F$ vanishes on $\ind_{(s,+\infty)}$ and on $\ind_{[s,+\infty)}$ for
every $s$. Deduce that $F$ vanishes on the indicator of every interval, then on every step function.
\item Deduce that $F(h) = 0$ for every $h \in C([\alpha,\beta])$. Where is item 7 used?
\item Apply this to $h = \cos$ and $h = \sin$ and conclude that $\mu = 0$.
\end{enumerate}

\medskip
\noindent\emph{Step 3: auditing the hypotheses.}
\begin{enumerate}[label=\textbf{\arabic*.}, resume]
\item Go back through your proof and say, for each of the three hypotheses on $\sigma$ --- bounded,
measurable, sigmoidal --- exactly where it was used. Is continuity of $\sigma$ used anywhere?
\end{enumerate}
\end{exercise}

% ================================================================
\section{Part E --- Assembling, and reading the result}

\begin{exercise}[Cybenko's theorem \tagcore]\label{ex:cybenko}
\leavevmode
\begin{enumerate}[label=\textbf{\arabic*.}]
\item Combine Exercises \ref{ex:skeleton} and \ref{ex:discr} into a single statement and write out
its proof in at most ten lines. This is the theorem of Cybenko (1989).
\item Your proof is by contradiction. Suppose a colleague hands you $f \in C([0,1]^2)$ and
$\varepsilon = 10^{-3}$ and asks for an explicit network. What does your proof give them?
\item Does the theorem provide any bound on the number $n$ of neurons as a function of
$\varepsilon$? Read your proof again and identify the precise step at which all quantitative
information is lost.
\item Does the theorem apply to $\sigma = \ReLU$? Justify your answer from the hypotheses, and say
what this suggests about how sharp those hypotheses are.
\end{enumerate}
\end{exercise}

\begin{exercise}[Reading the theorem critically \tagcore]\label{ex:reading}
\leavevmode
\begin{enumerate}[label=\textbf{\arabic*.}]
\item The theorem concerns $\Sigma(\sigma) = \bigcup_n \Sigma_n(\sigma)$. Give a heuristic
dimension-counting argument suggesting that $\Sigma_n(\sigma)$, for \emph{fixed} $n$, cannot be
dense in $C(K)$. Why is turning this into a proof delicate? \emph{(Recall Remark 1.7 of the lecture
notes.)}
\item A PINN will be trained by minimising a functional over $\theta$. List the three things a
practitioner would want to know that Cybenko's theorem does not address.
\item In one sentence each, say which later session takes charge of each of the three.
\end{enumerate}
\end{exercise}

% ================================================================
\section{Part F --- The complete characterisation}

\begin{setting}
\textbf{(SW)} \emph{Stone--Weierstrass.} A subalgebra of $C(K)$ that contains the constants and
separates the points of $K$ is dense in $C(K)$.

\smallskip
We shall prove the sufficiency half of the Leshno--Lin--Pinkus--Schocken theorem in the case
$\sigma \in C^\infty(\R)$. Throughout the exercise, $V := \overline{\Sigma(\sigma)}$, the closure
being taken in $C(K)$; note that $V$ is a closed linear subspace, so it is stable under uniform
limits --- this is the only property of $V$ that will be used, and it will be used constantly.
\end{setting}

\begin{exercise}[LLPS, smooth case \tagcore]\label{ex:llps}
Let $\sigma \in C^\infty(\R)$ be non-polynomial and $K \subset \R^d$ compact.

\medskip
\noindent\emph{Step 1: differentiating inside the class.}
\begin{enumerate}[label=\textbf{\arabic*.}]
\item Fix $w \in \R^d$, $b \in \R$, $\lambda \in \R$ and $h \ne 0$. Show that
\[
x \longmapsto \frac{\sigma\big((\lambda+h)\,w\cdot x + b\big) - \sigma\big(\lambda\,w\cdot x + b\big)}{h}
\]
belongs to $\Sigma_2(\sigma)$, hence to $V$.
\item Show that as $h \to 0$ this converges \emph{uniformly on $K$} to
$x \mapsto (w\cdot x)\,\sigma'(\lambda\,w\cdot x + b)$. \emph{Uniform convergence is the whole
point; say precisely why compactness of $K$ gives it.} Deduce that this limit lies in $V$.
\item Iterate to obtain, for every $k \in \N$,
\[
x \longmapsto (w\cdot x)^k\, \sigma^{(k)}\big(\lambda\, w\cdot x + b\big) \ \in\ V .
\]
\item Show that if $\sigma \in C^\infty$ is non-polynomial then for every $k$ there exists
$b_k \in \R$ with $\sigma^{(k)}(b_k) \ne 0$. \emph{(Argue by contraposition.)}
\item Deduce that $(w\cdot x)^k \in V$ for every $k \in \N$ and every $w \in \R^d$.
\end{enumerate}

\medskip
\noindent\emph{Step 2: from ridge monomials to all polynomials.}
\begin{enumerate}[label=\textbf{\arabic*.}, resume]
\item Prove the polarisation identity: for $y_1,\dots,y_k$ in a commutative ring,
\[
k!\, y_1 y_2\cdots y_k = \sum_{S \subseteq \{1,\dots,k\}} (-1)^{k-\abs S}\Big(\sum_{j\in S} y_j\Big)^{k}.
\]
\emph{(Expand by the multinomial formula and show that a monomial in which some $y_j$ is absent
contributes $0$ after summation over $S$.)} Check it by hand for $k = 2$ and $k=3$.
\item Deduce that every monomial $x_{i_1}\cdots x_{i_k}$ belongs to $V$, and then that
$\Pp(\R^d)\subseteq V$.
\item Conclude using \textup{(SW)}. Verify each hypothesis of \textup{(SW)} explicitly: subalgebra,
constants, separation of points.
\end{enumerate}

\medskip
\noindent\emph{Step 3: what the smoothness bought.}
\begin{enumerate}[label=\textbf{\arabic*.}, resume]
\item Where did you use $\sigma \in C^\infty$ rather than merely $\sigma \in C^0$? Sketch, in two or
three sentences, what one would have to do to treat a continuous $\sigma$.
\item Combine with Exercise \ref{ex:poly} to state the full theorem: a necessary and sufficient
condition on $\sigma \in C(\R)$ for $\Sigma(\sigma)$ to be dense in $C(K)$ for every compact $K$ and
every $d$.
\end{enumerate}
\end{exercise}

% ================================================================
\section{Part G --- Testing the hypotheses}

\begin{exercise}[What each hypothesis is doing \tagextra]\label{ex:hypotheses}
\leavevmode
\begin{enumerate}[label=\textbf{\arabic*.}]
\item \emph{Non-polynomiality is global, not local.} Does there exist $\sigma \in C^\infty(\R)$,
non-polynomial, with $\sigma^{(k)}(0) = 0$ for every $k \ge 3$? Exhibit one. Does the proof of
Exercise \ref{ex:llps} survive? Which step would have broken if we had insisted on taking
$\lambda = b = 0$ throughout?
\item \emph{Interior matters.} Let $K = \{0\} \cup \{e_1\} \subset \R^d$, two points. Identify
$C(K)$, and determine whether $\Sigma(\sigma)$ is dense in $C(K)$ when $\sigma(t) = t^2$. Compare
with Exercise \ref{ex:poly}, item 2, and explain the apparent conflict.
\item \emph{Compactness matters.} In Exercise \ref{ex:discr}, item 6, we used that $K$ is compact.
What breaks in the argument if $K = \R^d$? Is the conclusion still true for, say, bounded continuous
functions on $\R^d$ with the sup norm?
\item \emph{Boundedness is not an extra hypothesis, for continuous $\sigma$.} Show that a continuous
$\sigma : \R \to \R$ admitting finite limits at $+\infty$ and $-\infty$ is bounded. Deduce that in
Cybenko's theorem the word ``bounded'' may be dropped when $\sigma$ is assumed continuous and
sigmoidal.
\item \emph{Where continuity really sits.} Exercise \ref{ex:discr} never used continuity of
$\sigma$. Exercise \ref{ex:skeleton} did. Explain in one sentence the division of labour, and what
would go wrong in Exercise \ref{ex:skeleton} for a discontinuous sigmoidal $\sigma$ such as
$\ind_{[0,\infty)}$.
\end{enumerate}
\end{exercise}

% ================================================================
\newpage
\section{Hints}

\noindent
\textbf{Exercise \ref{ex:warmup}.}
\emph{3.} Treat $x > b$, $x \in (a,b)$, $x < a$ separately and use the two limits defining
sigmoidality. \emph{4.} Step functions are dense in $C([0,1])$ for the uniform norm --- or invoke
uniform continuity directly. \emph{5.} In $d \ge 2$ a ridge function is constant along a whole
hyperplane; what would the analogue of a ``bump'' be, and can you build a bump on a bounded region
out of finitely many of them?

\medskip\noindent
\textbf{Exercise \ref{ex:poly}.}
\emph{2.} A function that is not polynomial on any neighbourhood: take $x \mapsto \abs{x_1 - a}$.
\emph{4.} $(wx+b)^2 = w^2x^2 + 2wbx + b^2$; now vary $w$ and $b$. \emph{6.} Build products, then
monomials, then sums.

\medskip\noindent
\textbf{Exercise \ref{ex:skeleton}.}
\emph{1.} For $v \in V$ and $t \ne 0$, $\norm{v + tx_0} = \abs t\,\norm{x_0 + v/t} \ge \abs t\,\delta$.
\emph{3.} Argue by contradiction: if $\overline{\Sigma(\sigma)} \ne C(K)$, produce $L$, then $\mu$,
then apply the hypothesis to the functions $\sigma(w\cdot{\,\cdot\,}+b)$, which lie in
$\Sigma(\sigma)$.

\medskip\noindent
\textbf{Exercise \ref{ex:discr}.}
\emph{1.} Feed $\lambda w$ and $\lambda b + \varphi$. \emph{4.} The limit obtained in item 3 does
not depend on $\varphi$, but the expression does --- that is what lets you vary $\varphi$.
\emph{8.} $\ind_{(s,t]} = \ind_{(s,\infty)} - \ind_{(t,\infty)}$. \emph{9.} A uniform limit of
functions on which a bounded linear functional vanishes.

\medskip\noindent
\textbf{Exercise \ref{ex:llps}.}
\emph{2.} $w\cdot x$ ranges over a bounded set; $\sigma'$ is uniformly continuous on compacts.
\emph{4.} If for some $k$ one had $\sigma^{(k)} \equiv 0$ on $\R$, what is $\sigma$?
\emph{6.} In the expansion of $(\sum_{j\in S}y_j)^k$, fix a monomial $y^\alpha$ and sum its
coefficient over all $S$ containing the support of $\alpha$; if the support is proper you are
summing $(-1)^{k-\abs S}$ over a nonempty binomial family.

\medskip\noindent
\textbf{Exercise \ref{ex:hypotheses}.}
\emph{1.} A function that is flat to infinite order at a point, glued to something non-polynomial.
\emph{2.} $C(K) \cong \R^2$; you only need to hit two values. \emph{4.} Continuity gives
boundedness on a compact; the limits give it outside.

% ================================================================

\ifsolutions
\newpage
\section{Solutions}

\subsection*{Exercise \ref{ex:warmup}}

\textbf{1.} $c$ scales the output (vertical amplitude and, if $c<0$, a reflection); $b$ translates
horizontally, the graph being centred at $x = -b/w$; $w$ controls the \emph{steepness}, the slope at
the centre being $cw\,\sigma'(0)$. Large $\abs w$ makes the transition sharp --- this is the
observation the whole of Part D turns into a proof.

\textbf{2.} Take $w = 0$, $b = t_0$: the function $x \mapsto \sigma(0\cdot x + t_0) = \sigma(t_0)$
is constant and nonzero, so $x \mapsto \gamma\,\sigma(t_0)^{-1}\sigma(0\cdot x + t_0) = \gamma$ for
any $\gamma\in\R$, and this is an element of $\Sigma_1(\sigma)$.

\textbf{3.} The left-hand side is $\sigma(w_1x+b_1) - \sigma(w_2x+b_2)$ with $w_1 = w_2 = \lambda$,
$b_1 = -\lambda a$, $b_2 = -\lambda b$: an element of $\Sigma_2(\sigma)$ with coefficients $\pm1$.
For the limit, if $x > b$ then both arguments tend to $+\infty$ and the difference tends to
$1 - 1 = 0$; if $a < x < b$ then $\lambda(x-a)\to+\infty$ and $\lambda(x-b)\to-\infty$, giving
$1 - 0 = 1$; if $x < a$ both tend to $-\infty$, giving $0 - 0 = 0$. At $x = a$ or $x=b$ the
behaviour depends on $\sigma$ at a fixed point and no limit is claimed.

\textbf{4.} Let $f \in C([0,1])$ and $\varepsilon > 0$. By uniform continuity there is a step
function $s = \sum_{i} \gamma_i \ind_{(a_i,b_i)}$ with $\norm{f - s}_\infty \le \varepsilon/2$ off a
finite set. By item 3 each $\ind_{(a_i,b_i)}$ is a pointwise limit of elements of
$\Sigma_2(\sigma)$, so a finite combination approximates $s$. The implicit classical result is the
density of step functions in $C([0,1])$ for the uniform norm.

\emph{Two honest caveats to note, since they explain the shape of the real proof.} The convergence
in item 3 is pointwise off $\{a,b\}$, not uniform --- near the jump the approximation is always
bad --- so this sketch needs care to become a proof even in $d=1$. And the number of neurons
produced is that of the step approximation, which is where a rate would come from if one pushed the
argument; Cybenko's proof abandons that possibility entirely.

\textbf{5.} In $d \ge 2$, $x \mapsto \sigma(w\cdot x + b)$ is constant on every hyperplane
$w\cdot x = \text{const}$, so it has no localisation transverse to $w$: the object built in item 3
is a \emph{slab}, infinite in the $d-1$ directions orthogonal to $w$, not a bump. One can intersect
slabs to localise, but intersection is not a linear operation, and $\Sigma(\sigma)$ is a linear
span: sums of slab functions do not produce an indicator of a box. Some genuinely different idea is
needed --- either the Fourier/ridge decomposition of Part D, or the differentiation trick of Part F.

\subsection*{Exercise \ref{ex:poly}}

\textbf{1.} If $\sigma(t) = \sum_{j=0}^m a_jt^j$ then
$\sigma(w\cdot x+b) = \sum_j a_j (w\cdot x + b)^j$, and $(w\cdot x+b)^j$ expands into a polynomial
in $x$ of total degree $\le j \le m$. A finite linear combination of such functions stays in
$\Pp_m(\R^d)$, which is a vector space; hence $\Sigma(\sigma)\subseteq\Pp_m(\R^d)$.

\textbf{2.} $\Pp_m(\R^d)$ is finite-dimensional, hence closed in $C(K)$, so
$\overline{\Sigma(\sigma)}\subseteq\Pp_m(\R^d)$. If $K$ has nonempty interior, pick $a$ interior and
consider $f(x) = \abs{x_1 - a_1}$, which is continuous on $K$ but coincides with no polynomial on
any neighbourhood of $a$ (a polynomial is $C^\infty$, $f$ is not differentiable at $x_1 = a_1$).
Hence $f \notin \Pp_m(\R^d)$ and $\Sigma(\sigma)$ is not dense.

\textbf{3.} Only in the last step: to say that a function on $K$ is not the restriction of a
polynomial one needs enough room. On a finite set, every function is a polynomial restriction, and
the conclusion fails --- see Exercise \ref{ex:hypotheses}, item 2.

\textbf{4.} $\sigma(wx+b) = (wx+b)^2 = w^2x^2 + 2wbx + b^2$. Taking $b = 0$, $w=1$ gives $x^2$;
$w = 0$, $b = 1$ gives the constant $1$; $w = b = 1$ gives $x^2+2x+1$, whence $x$ by subtraction.
So $\Sigma(\sigma) \supseteq \Pp_2(\R)$, and the reverse inclusion is item 1: the inclusion is an
\emph{equality}, $\Sigma(t^2) = \Pp_2(\R)$.

\textbf{5.} By induction on the layers: an affine map sends polynomials of degree $\le D$ to
polynomials of degree $\le D$, and $t \mapsto t^2$ doubles the degree. Starting from degree $1$
(the first affine map), after $L-1$ hidden layers the degree is at most $2^{L-1}$. So
$\Nn^{t^2}_{L,W}(d)\subseteq \Pp_{2^{L-1}}(\R^d)$ for every width $W$, and this is a
finite-dimensional, hence closed, proper subspace of $C(K)$ when $K$ has interior: not dense.

\textbf{6.} The identity $xy = \tfrac14((x+y)^2 - (x-y)^2)$ says that one hidden layer of two
quadratic units followed by a linear combination computes the product of two quantities available at
that layer. Starting from the inputs $x_1,\dots,x_d$ and the constant $1$ (available as
$\sigma(0\cdot x + 1)$), one layer produces all degree-two monomials, $m$ layers of pairwise
products produce any monomial of degree $\le 2^m$, and the output layer sums. Hence
$\bigcup_{L,W}\Nn^{t^2}_{L,W}(d)\supseteq \Pp(\R^d)$, dense in $C(K)$ by Stone--Weierstrass.

\textbf{7.} No conflict: density fails at each \emph{fixed} depth and holds only for the union over
all depths. LLPS quantifies over $\Sigma(\sigma)$, the span of a \emph{single} hidden layer, and for
that class non-polynomiality is necessary and sufficient. Items 5--6 show that the analogous
statement for the union over all depths is false --- $t^2$ fails the LLPS criterion yet deep
quadratic networks are dense. Keeping track of the class a theorem quantifies over is not pedantry
here: it changes the truth value.

\subsection*{Exercise \ref{ex:skeleton}}

\textbf{1.} Let $x_0 \in X\setminus V$. Since $V$ is closed, $\delta := \dist(x_0,V) > 0$. On
$W := V \oplus \R x_0$ define $L(v + tx_0) := t\delta$; this is well defined because the sum is
direct ($x_0\notin V$), and linear. For $t \ne 0$,
\[
\norm{v + tx_0} = \abs t \cdot \norm{x_0 + v/t} \ \ge\ \abs t\,\delta = \abs{L(v+tx_0)},
\]
since $-v/t \in V$; for $t = 0$ the inequality is trivial. So $\norm{L}_{W'}\le 1$ and $L$ is
continuous. By (HB) extend it to $X$ with the same norm. Then $L|_V = 0$ and $L(x_0) = \delta \ne 0$,
so $L \ne 0$.

\textbf{2.} Closedness is used exactly once, to get $\delta > 0$: if $V$ were merely a subspace with
$\overline V = X$, every $x_0$ would be at distance $0$ from $V$ and the construction would produce
$L = 0$. Conceptually: a dense subspace is annihilated only by the zero functional, which is the
whole point.

\textbf{3.} $S := \Sigma(\sigma)$ is a linear subspace of $C(K)$ --- note this uses that
$x\mapsto\sigma(w\cdot x+b)$ is \emph{continuous} on $K$, so that $S \subseteq C(K)$ at all.
Suppose $\overline S \ne C(K)$. Then $\overline S$ is a closed proper subspace, so by item 1 there
is $L \in C(K)'$, $L \ne 0$, vanishing on $\overline S$. By (RM), $L(f) = \int_K f\dd\mu$ for a
finite signed regular Borel $\mu$, and $\mu \ne 0$ since $L\ne0$. For all $w,b$ the function
$\sigma(w\cdot{\,\cdot\,}+b)$ lies in $S$, so $\int_K\sigma(w\cdot x+b)\dd\mu(x) = 0$. As $\sigma$
is discriminatory on $K$, $\mu = 0$ --- a contradiction. Hence $\overline S = C(K)$.

\textbf{4.} At the very first line: without continuity, $\sigma(w\cdot{\,\cdot\,}+b)$ need not be an
element of $C(K)$, and neither $S \subseteq C(K)$ nor the application of $L$ to it makes sense. This
is the only place. Everything else --- including Part D --- works for measurable $\sigma$.

\textbf{5.} One would first need $K\subseteq[0,1]^d$ (a translation and dilation, harmless), and
then to extend the measure $\mu$, which lives on $K$, to a measure $\tilde\mu$ on $[0,1]^d$ by
$\tilde\mu(A) := \mu(A\cap K)$. This $\tilde\mu$ is a finite signed regular Borel measure on the
cube, it is nonzero if $\mu$ is, and
$\int_{[0,1]^d}\sigma(w\cdot x+b)\dd\tilde\mu = \int_K\sigma(w\cdot x+b)\dd\mu = 0$, so the cube
version applies and gives $\tilde\mu = 0$, hence $\mu = 0$.

The step is harmless but it \emph{is} a step, and it is routinely omitted. Stating
Definition \ref{def:discr} on an arbitrary compact costs nothing --- the proof of Exercise
\ref{ex:discr} uses only compactness --- and removes the need for it entirely.

\subsection*{Exercise \ref{ex:discr}}

\textbf{1.} Apply \eqref{eq:hyp} with $w' := \lambda w$ and $b' := \lambda b + \varphi$; then
$\sigma(w'\cdot x + b') = \sigma(\lambda(w\cdot x+b)+\varphi) = \sigma_\lambda(x)$.

\textbf{2.} If $w\cdot x + b > 0$ then $\lambda(w\cdot x+b)+\varphi \to +\infty$ and
$\sigma_\lambda(x)\to 1$; if $w\cdot x + b< 0$ it tends to $-\infty$ and $\sigma_\lambda(x)\to 0$;
if $w\cdot x + b = 0$ then $\sigma_\lambda(x) = \sigma(\varphi)$ for \emph{every} $\lambda$, so the
limit is $\sigma(\varphi)$. The third case is the one that carries the information.

\textbf{3.} $\abs{\sigma_\lambda}\le\norm\sigma_\infty$ uniformly in $\lambda$ and $x$, and the
constant $\norm\sigma_\infty$ is $\abs\mu$-integrable because $\abs\mu(K)<\infty$. Dominated
convergence for the finite signed measure $\mu$ (apply it to the Jordan decomposition
$\mu = \mu^+-\mu^-$) gives the exchange.

\textbf{4.} Passing to the limit in $\int\sigma_\lambda\dd\mu = 0$:
\[
0 = \int_K \Big(\ind_{H^+_{w,b}} + \sigma(\varphi)\ind_{\Pi_{w,b}}\Big)\dd\mu
= \mu(H^+_{w,b}) + \sigma(\varphi)\mu(\Pi_{w,b}).
\]
This holds for every $\varphi$. Letting $\varphi\to-\infty$ gives $\sigma(\varphi)\to0$, hence
$\mu(H^+_{w,b}) = 0$; letting $\varphi\to+\infty$ gives $\sigma(\varphi)\to1$, hence
$\mu(H^+_{w,b})+\mu(\Pi_{w,b}) = 0$, so $\mu(\Pi_{w,b}) = 0$ too. Note that the limit function does
not depend on $\varphi$ in the first two cases, but the value on $\Pi$ does: this is exactly what
makes the two limits informative.

\textbf{5.} $x\mapsto w\cdot x$ is continuous and $K$ compact, so the image is compact, hence
contained in $[\alpha,\beta]$ with $\alpha := \min_K w\cdot x$ and $\beta := \max_K w\cdot x$. It
need not \emph{equal} $[\alpha,\beta]$ because $K$ may be disconnected: for $K = \{0\}\cup\{1\}
\subset\R$ and $w = 1$, the image is $\{0,1\}$, not $[0,1]$. Nothing in the sequel needs equality,
only that a compact interval containing the image exists.

\textbf{6.} Linearity is clear. For the bound,
$\abs{F(h)} = \abs{\int_K h(w\cdot x)\dd\mu(x)} \le \int_K \abs{h(w\cdot x)}\dd\abs\mu(x)
\le \norm h_\infty\,\abs\mu(K)$.

\textbf{7.} $\{x \in K : w\cdot x > s\} = H^+_{w,-s}$ and
$\{x\in K : w\cdot x \ge s\} = H^+_{w,-s}\cup\Pi_{w,-s}$, so
$F(\ind_{(s,\infty)}) = \mu(H^+_{w,-s}) = 0$ and
$F(\ind_{[s,\infty)}) = \mu(H^+_{w,-s})+\mu(\Pi_{w,-s}) = 0$ by Step 1. For an interval,
$\ind_{(s,t]} = \ind_{(s,\infty)} - \ind_{(t,\infty)}$ and similarly for the other three types, so
$F$ vanishes on every interval indicator by linearity, hence on every step function.

\textbf{8.} Let $h\in C([\alpha,\beta])$. By uniform continuity there are step functions $s_n$ with
$\norm{h - s_n}_\infty\to0$. Then $\abs{F(h)} = \abs{F(h-s_n)}\le\norm{h-s_n}_\infty\abs\mu(K)\to0$,
so $F(h) = 0$. Item 6 is used precisely here: without the bound $\abs{F(h)}\le\norm h_\infty\abs\mu(K)$
one could not pass from step functions to their uniform limits.

\textbf{9.} Taking $h = \cos$ and $h=\sin$ restricted to $[\alpha,\beta]$, which are continuous,
\[
\widehat\mu(w) := \int_K \ee^{\ii w\cdot x}\dd\mu(x)
= \int_K\cos(w\cdot x)\dd\mu(x) + \ii\int_K\sin(w\cdot x)\dd\mu(x) = F(\cos)+\ii F(\sin) = 0 .
\]
This holds for every $w\in\R^d$ (each $w$ has its own $[\alpha,\beta]$, which is harmless). By the
uniqueness theorem for the Fourier--Stieltjes transform, $\mu = 0$.

\textbf{10.} \emph{Bounded} is used in item 3, for the domination. \emph{Measurable} is used
wherever $\sigma_\lambda$ is integrated, i.e.\ in items 1 and 3. \emph{Sigmoidal} is used twice: in
item 2 to identify the pointwise limit, and in item 4 to let $\sigma(\varphi)$ sweep from $0$ to
$1$ --- it is the existence of \emph{two distinct} limits that makes the two passages informative;
had $\sigma$ the same limit at both ends, item 4 would give a single equation instead of two.
Continuity of $\sigma$ is used \emph{nowhere} in this exercise.

\subsection*{Exercise \ref{ex:cybenko}}

\textbf{1.} \emph{Theorem (Cybenko, 1989).} Let $\sigma\in C(\R)$ be bounded and sigmoidal and let
$K\subset\R^d$ be compact. Then $\Sigma(\sigma)$ is dense in $C(K)$.

\emph{Proof.} $S := \Sigma(\sigma)$ is a linear subspace of $C(K)$, the functions
$\sigma(w\cdot{\,\cdot\,}+b)$ being continuous. Suppose $\overline S\ne C(K)$. By Exercise
\ref{ex:skeleton}, item 1, there is $L\in C(K)'$, $L\ne0$, with $L|_{\overline S} = 0$. By (RM)
there is a finite signed regular Borel measure $\mu\ne0$ on $K$ with $L(f) = \int_K f\dd\mu$. Since
$\sigma(w\cdot{\,\cdot\,}+b)\in S$ for all $w,b$, we get $\int_K\sigma(w\cdot x+b)\dd\mu(x) = 0$ for
all $w,b$. By Exercise \ref{ex:discr}, $\sigma$ is discriminatory on $K$, so $\mu = 0$ --- a
contradiction. Hence $\overline S = C(K)$. $\square$

\textbf{2.} Nothing. The proof shows that the assumption ``not dense'' is contradictory; it never
constructs an approximant. Handed $f$ and $\varepsilon$, one knows a network exists and has no
procedure to find one, nor even to tell how large it must be.

\textbf{3.} No bound. The quantitative information is destroyed at the passage through Hahn--Banach:
from that point on the argument manipulates an abstract functional and a measure, and $n$ has
disappeared from the discussion entirely --- indeed $n$ never appears in the proof at all, since we
work with $\Sigma(\sigma) = \bigcup_n\Sigma_n(\sigma)$ from the first line. Recovering a rate
requires a completely different argument: this is Barron's theorem, §\,3 of the lecture notes.

\textbf{4.} No: $\ReLU$ is unbounded, so it is not sigmoidal in the sense used here, and both
hypotheses of the theorem fail. This shows the hypotheses are \emph{not} sharp --- $\ReLU$ does
generate a dense class, as LLPS shows, since it is continuous and not a polynomial. Cybenko's
hypotheses are those his method needs (boundedness for the domination, sigmoidality for the step
limit), not those the conclusion needs. This is a healthy thing for students to see: a theorem's
hypotheses often describe its proof rather than its truth.

\subsection*{Exercise \ref{ex:reading}}

\textbf{1.} $\Sigma_n(\sigma)$ is the image of $\R^{n(d+2)}$ under the map
$(c_i,w_i,b_i)\mapsto\sum_i c_i\sigma(w_i\cdot{\,\cdot\,}+b_i)$: a set parametrised by finitely many
numbers inside the infinite-dimensional space $C(K)$. One expects such a set to be ``thin'', and in
particular not dense.

Turning this into a proof is delicate because the parameter space is \emph{unbounded}: the image of
a non-compact set under a continuous map need not be closed, and indeed Remark 1.7 of the lecture
notes records that $\Sigma_n(\sigma)$ is in general \emph{not} closed. So one cannot argue ``finite
dimensional, hence closed, hence not dense'' as in Exercise \ref{ex:poly}. A correct treatment needs
genuine work; see the references given in that remark.

\textbf{2.} (i) How many neurons are needed for accuracy $\varepsilon$, and how does that number
depend on $d$? (ii) Can the good parameters be \emph{found} by an optimisation algorithm? (iii) Is
it enough to approximate $u$, given that a PINN must approximate $\Dd[u]$, i.e.\ derivatives?

\textbf{3.} (i) is §\,3 and §\,4 of the lecture notes (Barron; Sobolev-norm rates), and is taken
further in Session 4. (ii) is Session 6, which analyses training and is where the practical gap
lies. (iii) is §\,4 of the lecture notes, and governs the choice of activation via Session 1,
Proposition 1.8.

\subsection*{Exercise \ref{ex:llps}}

\textbf{1.} The two terms are $\sigma(w'\cdot x + b)$ with $w' = (\lambda+h)w$ and $w'' = \lambda w$
respectively, so each lies in $\Sigma_1(\sigma)$; their difference divided by $h$ lies in
$\Sigma_2(\sigma)\subseteq V$.

\textbf{2.} Set $u := w\cdot x$, which ranges over a compact interval $I$ as $x$ ranges over $K$.
For fixed $b$, the map $(\lambda,u)\mapsto \sigma(\lambda u + b)$ is $C^\infty$; by the mean value
theorem,
\[
\frac{\sigma((\lambda+h)u+b)-\sigma(\lambda u+b)}{h} = u\,\sigma'(\theta_{h,u}\,u + b)
\]
for some $\theta_{h,u}$ between $\lambda$ and $\lambda+h$. As $\abs u$ is bounded on $I$ and
$\sigma'$ is uniformly continuous on the compact set $\{\theta u + b : \abs{\theta-\lambda}\le1,\ u \in I\}$,
the right-hand side converges to $u\,\sigma'(\lambda u + b)$ uniformly in $u \in I$, hence uniformly
in $x \in K$. Compactness of $K$ is what makes $I$ compact and the convergence uniform rather than
merely pointwise. Since $V$ is closed for the uniform norm, the limit is in $V$.

\textbf{3.} Repeat item 2 with $\sigma$ replaced by $\sigma'$, then $\sigma''$, and so on: each
differentiation in $\lambda$ multiplies by a factor $(w\cdot x)$ and raises the order of the
derivative by one. An induction on $k$ gives
$(w\cdot x)^k\sigma^{(k)}(\lambda w\cdot x + b)\in V$; all the intermediate functions are in $V$ and
$V$ is closed, so nothing escapes.

\textbf{4.} Contraposition. Suppose that for some $k$ we have $\sigma^{(k)}(b) = 0$ for every
$b\in\R$, i.e.\ $\sigma^{(k)}\equiv0$. Integrating $k$ times shows $\sigma$ is a polynomial of
degree $< k$, contradicting the hypothesis. Hence for every $k$ there is $b_k$ with
$\sigma^{(k)}(b_k)\ne0$. \emph{Note the quantifiers: non-polynomiality is a statement about
$\sigma^{(k)}$ on all of $\R$, not at a point.}

\textbf{5.} Take $\lambda = 0$ and $b = b_k$ in item 3: the function
$x\mapsto \sigma^{(k)}(b_k)\,(w\cdot x)^k$ lies in $V$, and $\sigma^{(k)}(b_k)\ne0$ is a nonzero
scalar, so $(w\cdot x)^k \in V$ by linearity.

\textbf{6.} Expand each $(\sum_{j\in S}y_j)^k$ by the multinomial formula. A monomial
$y^\alpha = y_1^{\alpha_1}\cdots y_k^{\alpha_k}$ with $\abs\alpha = k$ appears in the term indexed
by $S$ if and only if $\mathrm{supp}(\alpha)\subseteq S$, always with the same multinomial
coefficient $\binom{k}{\alpha}$. Its total coefficient in the sum is therefore
\[
\binom{k}{\alpha}\sum_{S \supseteq \mathrm{supp}(\alpha)}(-1)^{k-\abs S}
= \binom{k}{\alpha}\sum_{j=0}^{k-p}\binom{k-p}{j}(-1)^{k-p-j}
= \binom{k}{\alpha}\,(1-1)^{k-p},
\]
where $p := \abs{\mathrm{supp}(\alpha)}$ and we summed over the $\binom{k-p}{j}$ ways of adjoining
$j$ further indices. This vanishes unless $p = k$, i.e.\ unless every $y_j$ occurs, which forces
$\alpha = (1,\dots,1)$ and leaves the coefficient $\binom{k}{1,\dots,1} = k!$. Hence the sum equals
$k!\,y_1\cdots y_k$.

Check for $k=2$: $(y_1+y_2)^2 - y_1^2 - y_2^2 = 2y_1y_2$. \checkmark\ For $k = 3$:
$(y_1+y_2+y_3)^3 - \sum_{i<j}(y_i+y_j)^3 + \sum_i y_i^3 = 6y_1y_2y_3$, as one checks by expanding.

\textbf{7.} Fix $i_1,\dots,i_k\in\{1,\dots,d\}$ and apply the identity with $y_j := x_{i_j}$. For
$S\subseteq\{1,\dots,k\}$ put $w_S := \sum_{j\in S}e_{i_j}\in\R^d$, so that
$\sum_{j\in S}y_j = w_S\cdot x$. The identity reads
\[
k!\;x_{i_1}\cdots x_{i_k} = \sum_{S}(-1)^{k-\abs S}\,(w_S\cdot x)^k,
\]
and every term on the right lies in $V$ by item 5. Since $V$ is a vector space,
$x_{i_1}\cdots x_{i_k}\in V$. Every monomial of every degree is of this form (constants being the
case $k=0$), so by linearity $\Pp(\R^d)\subseteq V$.

\textbf{8.} $\Pp(\R^d)$ restricted to $K$ is a subalgebra of $C(K)$: it is a vector space, stable
under products. It contains the constants. It separates points: if $x\ne y$ in $K$, some coordinate
differs, say $x_i\ne y_i$, and the polynomial $z\mapsto z_i$ separates them. By (SW), $\Pp(\R^d)$ is
dense in $C(K)$. As $\Pp(\R^d)\subseteq V$ and $V$ is closed, $V = C(K)$, i.e.\ $\Sigma(\sigma)$ is
dense.

\textbf{9.} Smoothness is used in items 3 and 4: item 3 iterates the differentiation $k$ times and
needs $\sigma^{(k)}$ to exist for all $k$; item 4 characterises non-polynomiality through the
$\sigma^{(k)}$. For continuous $\sigma$ one regularises: convolve with a mollifier
$\varphi\in C_c^\infty$ to get $\sigma*\varphi\in C^\infty$, observe that
$\Sigma(\sigma*\varphi)\subseteq\overline{\Sigma(\sigma)}$ (a convolution is a limit of Riemann sums
of translates, and translates of $\sigma$ stay in $\Sigma(\sigma)$ once composed with $w\cdot x+b$),
and check that if $\sigma$ is not a polynomial then $\sigma*\varphi$ is not a polynomial for some
$\varphi$. Then apply the smooth case. The details are in Pinkus (1999), §\,3; they are more
delicate than this sketch suggests, which is why we did not do them.

\textbf{10.} \emph{Theorem (Leshno, Lin, Pinkus, Schocken, 1993).} Let $\sigma\in C(\R)$. Then
$\Sigma(\sigma)$ is dense in $C(K)$ for every compact $K\subset\R^d$ and every $d\ge1$ if and only
if $\sigma$ is not a polynomial. The ``only if'' is Exercise \ref{ex:poly}; the ``if'' is this
exercise for $\sigma\in C^\infty$, and item 9 for the general case.

\subsection*{Exercise \ref{ex:hypotheses}}

\textbf{1.} Yes: take
\[
\sigma(t) := \begin{cases} \ee^{-1/t^2} & t > 0,\\ 0 & t\le 0.\end{cases}
\]
This is $C^\infty$, all its derivatives vanish at $0$ (in particular for $k\ge3$), and it is not a
polynomial since it vanishes on a half-line without being identically zero.

The proof of Exercise \ref{ex:llps} survives untouched: item 4 asserts the existence, for each $k$,
of \emph{some} $b_k$ with $\sigma^{(k)}(b_k)\ne0$, and here one may take any $b_k>0$. What would
have broken is a version of item 5 insisting on $\lambda = b = 0$: that would require
$\sigma^{(k)}(0)\ne0$, which fails here for every $k\ge3$. The moral: non-polynomiality is a
\emph{global} property, and the proof must be allowed to look for a good point rather than being
handed one.

\textbf{2.} $K$ has two points, so $C(K)\cong\R^2$ via $f\mapsto(f(0),f(e_1))$, with the sup norm.
With $\sigma(t) = t^2$: taking $w = 0$, $b=1$ gives the vector $(1,1)$; taking $w = e_1$, $b = 0$
gives $(\sigma(0),\sigma(1)) = (0,1)$. These two vectors are independent, so
$\Sigma_2(\sigma) = \R^2 = C(K)$: $\Sigma(\sigma)$ is not merely dense, it is \emph{everything}.

There is no conflict with Exercise \ref{ex:poly}, item 2, whose conclusion was conditional on $K$
having nonempty interior. The inclusion $\Sigma(\sigma)\subseteq\Pp_2(\R^d)$ still holds; what fails
is that $\Pp_2(\R^d)$ is a proper subspace \emph{of $C(K)$}, because restriction to a two-point set
is far from injective on polynomials. This is exactly the role of the interior hypothesis, and it is
worth having seen it fail once.

\textbf{3.} Two things break. First, $\{w\cdot x : x\in K\}$ is no longer contained in a compact
interval, so the step-function approximation of item 8 of Exercise \ref{ex:discr} --- which relies on
uniform approximation of a continuous function on a \emph{compact} interval --- is unavailable.
Second, and more fundamentally, $C_b(\R^d)$ with the sup norm has a much larger dual than the space
of finite measures, so (RM) in the form we used does not apply and the skeleton of Exercise
\ref{ex:skeleton} collapses.

The conclusion is in fact false in that setting for a simple reason: uniform approximation on all of
$\R^d$ forces matching behaviour at infinity, and every element of $\Sigma(\sigma)$ for sigmoidal
$\sigma$ has strong constraints there, whereas a general bounded continuous function has none. The
right statement is the one we proved: density on \emph{every} compact, that is, uniform convergence
on compacts.

\textbf{4.} Let $\ell_\pm := \lim_{t\to\pm\infty}\sigma(t)$ be finite. There is $T>0$ such that
$\abs{\sigma(t) - \ell_+}\le1$ for $t\ge T$ and $\abs{\sigma(t)-\ell_-}\le1$ for $t\le-T$, so
$\abs\sigma\le\max(\abs{\ell_+},\abs{\ell_-})+1$ outside $[-T,T]$. On the compact $[-T,T]$ the
continuous function $\sigma$ is bounded. Hence $\sigma$ is bounded on $\R$.

Consequently, in the statement of Cybenko's theorem, ``bounded'' is redundant once $\sigma$ is
assumed continuous and sigmoidal: it is implied. It is \emph{not} redundant in Exercise
\ref{ex:discr}, where $\sigma$ is only assumed measurable.

\textbf{5.} The division of labour is clean. Exercise \ref{ex:discr} is a measure-theoretic
statement about the family $\{\sigma(w\cdot{\,\cdot\,}+b)\}$ tested against $\mu$; integration needs
only measurability and boundedness. Exercise \ref{ex:skeleton} is a functional-analytic statement
inside $C(K)$; it needs those functions to \emph{be} elements of $C(K)$, hence continuity.

For $\sigma = \ind_{[0,\infty)}$: Exercise \ref{ex:discr} applies verbatim and $\sigma$ is
discriminatory. But $\Sigma(\sigma)$ consists of step-like functions, which are not continuous, so
$\Sigma(\sigma)\not\subseteq C(K)$ and the question ``is $\Sigma(\sigma)$ dense in $C(K)$?'' is not
even well posed as stated. One can of course ask about density in $L^p(K)$ instead, and there the
argument goes through with the corresponding duality --- which is how Cybenko's result is usually
extended.

% ================================================================
\newpage
\fi

\section*{References}
\addcontentsline{toc}{section}{References}

\begin{itemize}[leftmargin=*, itemsep=4pt]
\item A.~Pinkus, \emph{Approximation theory of the MLP model in neural networks}, Acta Numerica
\textbf{8} (1999), 143--195. --- \emph{The reference for this whole sheet; §\,3 contains the
reduction from continuous to smooth $\sigma$ sketched in Exercise \ref{ex:llps}, item 9.}
\item G.~Cybenko, \emph{Approximation by superpositions of a sigmoidal function}, Mathematics of
Control, Signals and Systems \textbf{2} (1989), 303--314.
\item M.~Leshno, V.~Y.~Lin, A.~Pinkus, S.~Schocken, \emph{Multilayer feedforward networks with a
nonpolynomial activation function can approximate any function}, Neural Networks \textbf{6} (1993),
861--867.
\item R.~DeVore, B.~Hanin, G.~Petrova, \emph{Neural network approximation}, Acta Numerica
\textbf{30} (2021), 327--444.
\end{itemize}

\end{document}
